% Application Behaviour Usage Guide
% A complete introduction to application behaviour with architecture diagrams and examples.
% Compile with: xelatex application_architecture.tex (run twice for TOC)
\documentclass[a4paper,11pt]{article}

% ============================================================
% Packages
% ============================================================
\usepackage{fontspec}
\usepackage{xeCJK}
\setCJKmainfont[BoldFont=STHeiti, ItalicFont=STKaiti]{STSong}
\setCJKsansfont{STHeiti}
\setCJKmonofont{STFangsong}

\usepackage[margin=2.5cm]{geometry}
\usepackage{graphicx}
\usepackage{xcolor}
\usepackage{hyperref}
\usepackage{listings}
\usepackage{booktabs}
\usepackage{longtable}
\usepackage{enumitem}
\usepackage{tikz}
\usepackage{tcolorbox}
\usepackage{fancyhdr}
\usepackage{titlesec}
\usepackage{amsmath}
\usepackage{amssymb}
\usepackage{float}
\usepackage{pifont}

\usetikzlibrary{arrows.meta, positioning, shapes.geometric, fit, calc, backgrounds, decorations.pathreplacing}

\tcbuselibrary{listings, skins, breakable}

% ============================================================
% Color Definitions
% ============================================================
\definecolor{erlred}{HTML}{A4070A}
\definecolor{erlblue}{HTML}{2B5EA7}
\definecolor{erlgreen}{HTML}{2E7D32}
\definecolor{erlyellow}{HTML}{F9A825}
\definecolor{erlpurple}{HTML}{6A1B9A}
\definecolor{erlmodule}{HTML}{D84315}
\definecolor{erlinit}{HTML}{00897B}
\definecolor{erlpid}{HTML}{E65100}
\definecolor{erlname}{HTML}{6A1B9A}
\definecolor{codebg}{HTML}{F5F5F5}
\definecolor{codeborder}{HTML}{CCCCCC}
\definecolor{tipbg}{HTML}{E8F5E9}
\definecolor{warnbg}{HTML}{FFF3E0}
\definecolor{notebg}{HTML}{E3F2FD}
\definecolor{darkgray}{HTML}{333333}

% ============================================================
% Hyperref Setup
% ============================================================
\hypersetup{
    colorlinks=true,
    linkcolor=erlblue,
    urlcolor=erlpurple,
    citecolor=erlgreen,
    pdftitle={Application Behaviour 用法指南},
    pdfauthor={Erlang OTP Tutorial},
    bookmarksnumbered=true,
}

% ============================================================
% Code Listing Style
% ============================================================
\lstdefinelanguage{erlang}{
    morekeywords={module, export, import, behaviour, behavior, spec, type,
        record, define, include, ifdef, ifndef, else, endif, undef,
        case, of, end, if, receive, after, when, fun, try, catch, throw,
        begin, spawn, spawn_link, register, self, send, ok, error,
        true, false, undefined, infinity, timeout, noreply, reply, stop,
        hibernate, ignore, normal, process_flag, trap_exit, monitor_node,
        application, start, start_link, start_phase, prep_stop, config_change},
    sensitive=true,
    morecomment=[l]{\%},
    morestring=[b]",
    morestring=[b]',
}

\lstdefinestyle{erlangstyle}{
    language=erlang,
    basicstyle=\ttfamily\small,
    keywordstyle=\color{erlblue}\bfseries,
    commentstyle=\color{erlgreen}\itshape,
    stringstyle=\color{erlred},
    numberstyle=\tiny\color{gray},
    numbers=left,
    numbersep=8pt,
    backgroundcolor=\color{codebg},
    frame=single,
    rulecolor=\color{codeborder},
    breaklines=true,
    breakatwhitespace=true,
    tabsize=4,
    showstringspaces=false,
    captionpos=b,
    xleftmargin=1.5em,
    framexleftmargin=1em,
    aboveskip=1em,
    belowskip=0.8em,
}

\lstdefinestyle{shellstyle}{
    basicstyle=\ttfamily\small\color{white},
    backgroundcolor=\color{darkgray},
    rulecolor=\color{darkgray},
    frame=single,
    breaklines=true,
    showstringspaces=false,
    xleftmargin=1.5em,
    framexleftmargin=1em,
    aboveskip=1em,
    belowskip=0.8em,
}

\lstset{style=erlangstyle}

% ============================================================
% Custom tcolorbox Environments
% ============================================================
\newtcolorbox{tipbox}[1][]{
    colback=tipbg, colframe=erlgreen!80!black,
    fonttitle=\bfseries, title={#1},
    arc=2mm, boxrule=0.8pt,
    left=6pt, right=6pt, top=4pt, bottom=4pt,
    breakable,
}

\newtcolorbox{warnbox}[1][]{
    colback=warnbg, colframe=erlyellow!80!black,
    fonttitle=\bfseries, title={#1},
    arc=2mm, boxrule=0.8pt,
    left=6pt, right=6pt, top=4pt, bottom=4pt,
    breakable,
}

\newtcolorbox{notebox}[1][]{
    colback=notebg, colframe=erlblue!80!black,
    fonttitle=\bfseries, title={#1},
    arc=2mm, boxrule=0.8pt,
    left=6pt, right=6pt, top=4pt, bottom=4pt,
    breakable,
}

% ============================================================
% Header / Footer
% ============================================================
\pagestyle{fancy}
\fancyhf{}
\fancyhead[L]{\small\textit{Application Behaviour 用法指南}}
\fancyhead[R]{\small\textit{Erlang/OTP}}
\fancyfoot[C]{\thepage}
\renewcommand{\headrulewidth}{0.4pt}

% ============================================================
% Title Formatting
% ============================================================
\titleformat{\section}
    {\LARGE\bfseries\color{erlblue}}
    {\thesection}{1em}{}[\vspace{-0.5em}\rule{\textwidth}{0.8pt}]

\titleformat{\subsection}
    {\Large\bfseries\color{erlblue!80!black}}
    {\thesubsection}{0.8em}{}

\titleformat{\subsubsection}
    {\large\bfseries\color{erlblue!60!black}}
    {\thesubsubsection}{0.6em}{}

% ============================================================
% Document Begin
% ============================================================
\begin{document}

% ---- Title Page ----
\begin{titlepage}
\centering
\vspace*{3cm}

{\fontsize{42}{48}\selectfont\bfseries\color{erlblue} application}\\[12pt]
{\fontsize{24}{28}\selectfont\color{erlblue!70!black} Behaviour 用法指南}

\vspace{2cm}

\begin{tikzpicture}
    \node[draw=erlblue, fill=erlblue!10, rounded corners=8pt, minimum width=4cm, minimum height=1.8cm, font=\large\bfseries] (ac) at (0,0) {App Controller};
    \node[draw=erlpurple, fill=erlpurple!10, rounded corners=8pt, minimum width=4cm, minimum height=1.8cm, font=\large\bfseries] (am) at (0,-3) {App Master};
    \node[draw=erlred, fill=erlred!10, rounded corners=8pt, minimum width=4cm, minimum height=1.8cm, font=\large\bfseries] (sup) at (0,-6) {Top Supervisor};

    \draw[-{Stealth[length=8pt]}, line width=2pt, erlblue]
        (ac.south) -- node[right, font=\small] {start\_application} (am.north);
    \draw[-{Stealth[length=8pt]}, line width=2pt, erlpurple]
        (am.south) -- node[right, font=\small] {Module:start/2} (sup.north);

    \node[right=15mm of ac, font=\small\color{gray}, text width=4cm] {每个节点一个\\管理所有 Application};
    \node[right=15mm of am, font=\small\color{gray}, text width=4cm] {每个 Application 一个\\Group Leader};
    \node[right=15mm of sup, font=\small\color{gray}, text width=4cm] {监督树根节点\\管理子进程};
\end{tikzpicture}

\vspace{2cm}

{\large\color{darkgray} 以 BSC（基站控制器）为例，从架构到实践}

\vspace{1cm}

{\normalsize\color{gray} Erlang/OTP $\cdot$ Application Behaviour $\cdot$ Supervision Tree}

\vfill
{\small\today}
\end{titlepage}

% ---- Table of Contents ----
\newpage
\tableofcontents
\newpage


% ============================================================
% Section 1: Overview
% ============================================================
\section{概述}

\subsection{什么是 Application？}

在 OTP 中，\textbf{Application} 是一个实现了特定功能的组件，它可以作为一个整体被启动和停止，并且可以在其他系统中被复用。Application 是 Erlang/OTP 系统中组织代码的最高层级单元。

\begin{tipbox}[核心思想]
\textbf{Application = Application Resource File + Callback Module + Supervision Tree}

\begin{itemize}[leftmargin=1.5em]
    \item \textbf{Application Resource File}（\texttt{.app} 文件）：描述 Application 的元数据（名称、版本、依赖、启动模块等）
    \item \textbf{Callback Module}（实现 \texttt{application} behaviour）：定义如何启动和停止 Application
    \item \textbf{Supervision Tree}：Application 内部的进程组织结构，由 Top Supervisor 管理
\end{itemize}
\end{tipbox}

\subsection{Application 的两种类型}

\begin{table}[H]
\centering
\begin{tabular}{@{}lll@{}}
\toprule
\textbf{类型} & \textbf{说明} & \textbf{示例} \\
\midrule
Active Application  & 有监督树，通过 \texttt{mod} 键指定回调模块 & kernel, stdlib, sasl \\
Library Application & 无监督树，仅提供工具模块 & crypto（部分） \\
\bottomrule
\end{tabular}
\caption{Application 的两种类型}
\end{table}

\subsection{Application Behaviour 提供的能力}

\begin{table}[H]
\centering
\begin{tabular}{@{}ll@{}}
\toprule
\textbf{能力} & \textbf{说明} \\
\midrule
统一启停管理       & \texttt{application:start/1,2} 和 \texttt{application:stop/1} \\
依赖管理           & 自动检查 \texttt{applications} 列表中的依赖是否已启动 \\
环境变量           & 通过 \texttt{application:get\_env/2,3} 和 \texttt{set\_env/3,4} 管理配置 \\
分布式支持         & 支持 failover 和 takeover \\
热代码升级         & 通过 \texttt{config\_change/3} 回调处理配置变更 \\
包含应用           & 通过 \texttt{included\_applications} 支持应用嵌套 \\
启动阶段           & 通过 \texttt{start\_phases} 支持分阶段启动 \\
\bottomrule
\end{tabular}
\caption{Application Behaviour 提供的能力}
\end{table}

\subsection{OTP 源码中的工作流}

以下直接对应 \texttt{application.erl} 源码中的回调定义：

\begin{tcolorbox}[colback=codebg, colframe=codeborder, arc=2mm, boxrule=0.5pt,
                  title={\textbf{application.erl --- Callback Functions}},
                  fonttitle=\bfseries\ttfamily, width=\textwidth]
\small\ttfamily
\begin{tabular}{@{}lcl@{}}
application module &  & Callback module \\
\textsf{-{}-{}-{}-{}-{}-{}-{}-{}-{}-{}-{}-{}-{}-{}-{}-{}-{}-} & & \textsf{-{}-{}-{}-{}-{}-{}-{}-{}-{}-{}-{}-{}-{}-{}-{}-{}-{}-{}-{}-{}-{}-} \\
application:start/1,2         & -{}-{}-{}-{}-> & Module:start/2 \\
application:stop/1            & -{}-{}-{}-{}-> & Module:prep\_stop/1 (optional) \\
                              & -{}-{}-{}-{}-> & Module:stop/1 \\
-                             & -{}-{}-{}-{}-> & Module:start\_phase/3 (optional) \\
-                             & -{}-{}-{}-{}-> & Module:config\_change/3 (optional) \\
\end{tabular}
\end{tcolorbox}


% ============================================================
% Section 2: Architecture Diagram
% ============================================================
\section{架构图}

Application 的启动涉及三个关键角色：Application Controller、Application Master 和 Top Supervisor。

\begin{figure}[H]
\centering
\includegraphics[width=\textwidth]{images/application_architecture.pdf}
\caption{Application Behaviour 架构图}
\label{fig:arch}
\end{figure}

\subsection{Application Controller}

Application Controller 是每个 Erlang 节点上的一个单例进程（注册名为 \texttt{application\_controller}），在节点启动时自动创建。它负责：

\begin{itemize}[leftmargin=2em]
    \item 加载和卸载 Application（\texttt{application:load/1}、\texttt{application:unload/1}）
    \item 启动和停止 Application（\texttt{application:start/1,2}、\texttt{application:stop/1}）
    \item 检查依赖关系（\texttt{applications} 列表）
    \item 管理 Application 的环境变量
\end{itemize}

\subsection{Application Master}

每个正在运行的 Application 都有一个 Application Master 进程。它由 Application Controller 创建，负责：

\begin{itemize}[leftmargin=2em]
    \item 调用回调模块的 \texttt{Module:start/2} 启动监督树
    \item 成为 Application 中所有进程的 Group Leader（用于标识进程归属）
    \item 在 Application 停止时，调用 \texttt{Module:prep\_stop/1} 和 \texttt{Module:stop/1}
    \item 监控 Top Supervisor，如果 Top Supervisor 终止，Application 也随之终止
\end{itemize}

\subsection{启动流程}

\begin{enumerate}[leftmargin=2em]
    \item \texttt{application:start(myapp)} --- 用户调用启动函数
    \item Application Controller 加载 \texttt{myapp.app} 文件（如果尚未加载）
    \item Application Controller 检查 \texttt{applications} 列表中的依赖是否已启动
    \item Application Controller 创建 Application Master
    \item Application Master 调用 \texttt{Module:start(normal, StartArgs)}
    \item \texttt{Module:start/2} 返回 \texttt{\{ok, Pid\}}，其中 \texttt{Pid} 是 Top Supervisor 的 PID
    \item Application 进入运行状态
\end{enumerate}

\subsection{停止流程}

\begin{enumerate}[leftmargin=2em]
    \item \texttt{application:stop(myapp)} --- 用户调用停止函数
    \item Application Master 调用 \texttt{Module:prep\_stop(State)}（如果定义了该回调）
    \item Application Master 通知 Top Supervisor 关闭
    \item 整个监督树按照启动的\textbf{逆序}终止
    \item Application Master 调用 \texttt{Module:stop(State)}
    \item Application Master 终止
    \item Application 仍处于已加载状态（可以重新启动）
\end{enumerate}


% ============================================================
% Section 3: Application Resource File
% ============================================================
\section{Application Resource File（.app 文件）}

每个 Application 都需要一个 \texttt{.app} 文件，通常位于 \texttt{ebin/} 目录下。

\subsection{基本格式}

\begin{lstlisting}[style=erlangstyle, caption={Application Resource File 基本格式}]
{application, AppName,
 [{description, Description},
  {vsn, Version},
  {modules, [Module1, Module2, ...]},
  {registered, [RegName1, RegName2, ...]},
  {applications, [kernel, stdlib, ...]},
  {env, [{Key1, Val1}, {Key2, Val2}]},
  {mod, {CallbackModule, StartArgs}}
 ]}.
\end{lstlisting}

\subsection{关键字段说明}

\begin{table}[H]
\centering
\begin{tabular}{@{}llp{7cm}@{}}
\toprule
\textbf{字段} & \textbf{类型} & \textbf{说明} \\
\midrule
\texttt{description} & string() & Application 的描述信息 \\
\texttt{vsn}         & string() & 版本号 \\
\texttt{modules}     & [atom()] & Application 包含的所有模块列表 \\
\texttt{registered}  & [atom()] & Application 注册的进程名列表 \\
\texttt{applications} & [atom()] & 依赖的 Application 列表（必须先启动） \\
\texttt{env}         & [{atom(), term()}] & 环境变量（配置参数） \\
\texttt{mod}         & \{Module, Args\} & 回调模块和启动参数 \\
\texttt{included\_applications} & [atom()] & 包含的子 Application \\
\texttt{start\_phases} & [{atom(), term()}] & 启动阶段定义 \\
\texttt{maxT}        & timeout() & Application 最大运行时间 \\
\texttt{maxP}        & integer() | infinity & 最大进程数 \\
\bottomrule
\end{tabular}
\caption{.app 文件关键字段}
\end{table}

\subsection{实例：BSC Application}

\begin{lstlisting}[style=erlangstyle, caption={bsc.app --- 基站控制器 Application 资源文件}]
{application, bsc,
 [{description, "Base Station Controller"},
  {vsn, "1.0"},
  {modules, [bsc,
             bsc_sup,
             frequency,
             freq_overload,
             logger,
             simple_phone_sup,
             phone_fsm
            ]},
  {registered, [bsc_sup, frequency, freq_overload,
                simple_phone_sup]},
  {applications, [kernel, stdlib, sasl]},
  {env, []},
  {mod, {bsc, []}}
 ]}.
\end{lstlisting}

\begin{notebox}[Library Application]
如果 Application 没有监督树（Library Application），则不需要 \texttt{mod} 字段。这类 Application 只提供工具模块，不启动任何进程。
\end{notebox}


% ============================================================
% Section 4: Callback Functions
% ============================================================
\section{回调函数}

\subsection{必须实现的回调}

\subsubsection{Module:start/2}

\begin{lstlisting}[style=erlangstyle]
-callback start(StartType, StartArgs) ->
    {ok, Pid} |
    {ok, Pid, State} |
    {error, Reason}.

-type start_type() :: normal
                    | {takeover, Node :: node()}
                    | {failover, Node :: node()}.
\end{lstlisting}

当 \texttt{application:start/1,2} 被调用时，Application Master 会调用此函数。它应该启动 Application 的 Top Supervisor 并返回其 PID。

\textbf{StartType 参数说明：}

\begin{table}[H]
\centering
\begin{tabular}{@{}lp{9cm}@{}}
\toprule
\textbf{StartType} & \textbf{说明} \\
\midrule
\texttt{normal} & 正常启动 \\
\texttt{\{takeover, Node\}} & 分布式 Application 从 Node 接管 \\
\texttt{\{failover, Node\}} & 分布式 Application 从 Node 故障转移 \\
\bottomrule
\end{tabular}
\caption{StartType 参数}
\end{table}

\subsubsection{Module:stop/1}

\begin{lstlisting}[style=erlangstyle]
-callback stop(State) -> term().
\end{lstlisting}

在 Application 停止后调用，用于清理工作。\texttt{State} 来自 \texttt{Module:prep\_stop/1} 的返回值（如果定义了该回调），否则来自 \texttt{Module:start/2} 的返回值。返回值被忽略。

\subsection{可选回调}

\subsubsection{Module:prep\_stop/1}

\begin{lstlisting}[style=erlangstyle]
-callback prep_stop(State) -> NewState.
\end{lstlisting}

在 Application 即将停止时、关闭进程之前调用。可以用于在进程终止前保存状态或执行清理操作。\texttt{NewState} 会传递给 \texttt{Module:stop/1}。

\subsubsection{Module:start\_phase/3}

\begin{lstlisting}[style=erlangstyle]
-callback start_phase(Phase, StartType, PhaseArgs) ->
    ok | {error, Reason}.
\end{lstlisting}

当 \texttt{.app} 文件中定义了 \texttt{start\_phases} 时，每个启动阶段都会调用此函数。用于在不同 Application 之间同步启动过程。

\subsubsection{Module:config\_change/3}

\begin{lstlisting}[style=erlangstyle]
-callback config_change(Changed, New, Removed) -> ok.
%% Changed :: [{Par, Val}]  - changed parameters
%% New     :: [{Par, Val}]  - new parameters
%% Removed :: [Par]         - removed parameters
\end{lstlisting}

在代码替换后，如果配置参数发生了变化，会调用此函数。


% ============================================================
% Section 5: Application API
% ============================================================
\section{常用 API}

\subsection{启动与停止}

\begin{table}[H]
\centering
\begin{tabular}{@{}lp{8cm}@{}}
\toprule
\textbf{函数} & \textbf{说明} \\
\midrule
\texttt{application:start(App)} & 以 \texttt{temporary} 类型启动 Application \\
\texttt{application:start(App, Type)} & 以指定类型启动 Application \\
\texttt{application:stop(App)} & 停止 Application \\
\texttt{application:ensure\_all\_started(App)} & 递归启动 Application 及其所有依赖 \\
\texttt{application:ensure\_started(App)} & 启动 Application，已启动则返回 \texttt{ok} \\
\bottomrule
\end{tabular}
\caption{启动与停止 API}
\end{table}

\subsection{restart\_type() --- 重启类型}

\begin{table}[H]
\centering
\begin{tabular}{@{}lp{9cm}@{}}
\toprule
\textbf{类型} & \textbf{说明} \\
\midrule
\texttt{permanent} & Application 终止时，所有其他 Application 和整个节点都会终止 \\
\texttt{transient} & 正常终止只报告；异常终止会导致所有其他 Application 和节点终止 \\
\texttt{temporary} & Application 终止只报告，不影响其他 Application（默认） \\
\bottomrule
\end{tabular}
\caption{restart\_type() 重启类型}
\end{table}

\begin{warnbox}[注意]
\texttt{transient} 类型在实际中用处有限，因为当监督树终止时，原因通常是 \texttt{shutdown} 而非 \texttt{normal}。
\end{warnbox}

\subsection{环境变量管理}

\begin{table}[H]
\centering
\begin{tabular}{@{}lp{7cm}@{}}
\toprule
\textbf{函数} & \textbf{说明} \\
\midrule
\texttt{application:get\_env(App, Key)} & 获取配置参数值 \\
\texttt{application:get\_env(App, Key, Default)} & 获取配置参数值，带默认值 \\
\texttt{application:get\_all\_env(App)} & 获取所有配置参数 \\
\texttt{application:set\_env(App, Key, Val)} & 设置配置参数 \\
\texttt{application:unset\_env(App, Key)} & 删除配置参数 \\
\bottomrule
\end{tabular}
\caption{环境变量管理 API}
\end{table}

\subsection{查询 API}

\begin{table}[H]
\centering
\begin{tabular}{@{}lp{7cm}@{}}
\toprule
\textbf{函数} & \textbf{说明} \\
\midrule
\texttt{application:which\_applications()} & 列出所有正在运行的 Application \\
\texttt{application:loaded\_applications()} & 列出所有已加载的 Application \\
\texttt{application:get\_application(Pid)} & 获取进程所属的 Application \\
\texttt{application:get\_supervisor(App)} & 获取 Application 的 Top Supervisor PID \\
\texttt{application:get\_key(App, Key)} & 获取 .app 文件中的配置键值 \\
\bottomrule
\end{tabular}
\caption{查询 API}
\end{table}


% ============================================================
% Section 6: Example — BSC Application
% ============================================================
\section{实例：BSC Application}

以一个基站控制器（Base Station Controller）为例，展示 Application Behaviour 的完整结构。

\subsection{回调模块：bsc.erl}

\begin{lstlisting}[style=erlangstyle, caption={bsc.erl --- Application 回调模块}]
-module(bsc).
-behaviour(application).

%% Application callbacks
-export([start/2, stop/1]).

start(_StartType, _StartArgs) ->
    bsc_sup:start_link().

stop(_Data) ->
    ok.
\end{lstlisting}

\begin{tipbox}[最小实现]
Application 回调模块的最小实现只需要两个函数：
\begin{itemize}[leftmargin=1.5em]
    \item \texttt{start/2}：启动 Top Supervisor，返回 \texttt{\{ok, Pid\}}
    \item \texttt{stop/1}：清理工作，通常返回 \texttt{ok}
\end{itemize}
\end{tipbox}

\subsection{Top Supervisor：bsc\_sup.erl}

\begin{lstlisting}[style=erlangstyle, caption={bsc\_sup.erl --- Top Supervisor}]
-module(bsc_sup).
-behaviour(supervisor).

-export([start_link/0, init/1]).
-export([stop/0]).

start_link() ->
    {ok, Pid} = supervisor:start_link(
        {local, ?MODULE}, ?MODULE, []),
    {ok, Pid}.

stop() ->
    exit(whereis(?MODULE), shutdown).

init(_) ->
    ChildSpecList = [child(freq_overload),
                     child(frequency),
                     child(simple_phone_sup)],
    {ok, {{rest_for_one, 2, 3600}, ChildSpecList}}.

child(Module) ->
    {Module, {Module, start_link, []},
     permanent, 2000, worker, [Module]}.
\end{lstlisting}

\subsection{Application Resource File：bsc.app}

\begin{lstlisting}[style=erlangstyle, caption={bsc.app --- Application 资源文件}]
{application, bsc,
 [{description, "Base Station Controller"},
  {vsn, "1.0"},
  {modules, [bsc, bsc_sup, frequency,
             freq_overload, logger,
             simple_phone_sup, phone_fsm]},
  {registered, [bsc_sup, frequency,
                freq_overload, simple_phone_sup]},
  {applications, [kernel, stdlib, sasl]},
  {env, []},
  {mod, {bsc, []}}
 ]}.
\end{lstlisting}


% ============================================================
% Section 7: Start Phases
% ============================================================
\section{启动阶段（Start Phases）}

当多个 Application 之间需要在启动过程中进行同步时，可以使用 Start Phases 机制。

\subsection{配置 Start Phases}

在 \texttt{.app} 文件中定义 \texttt{start\_phases}：

\begin{lstlisting}[style=erlangstyle, caption={带 start\_phases 的 .app 文件}]
{application, bsc,
 [{description, "Base Station Controller"},
  {vsn, "1.0"},
  {modules, [bsc, bsc_sup, ...]},
  {registered, [bsc_sup, ...]},
  {applications, [kernel, stdlib, sasl]},
  {start_phases, [{init, []},
                  {admin, []},
                  {oper, []}]},
  {env, []},
  {mod, {bsc, []}}
 ]}.
\end{lstlisting}

\subsection{实现 start\_phase/3 回调}

\begin{lstlisting}[style=erlangstyle, caption={带 start\_phase 回调的 bsc.erl}]
-module(bsc).
-behaviour(application).

-export([start/2, start_phase/3, stop/1]).

start(_StartType, _StartArgs) ->
    bsc_sup:start_link().

start_phase(StartPhase, StartType, Args) ->
    io:format("bsc:start_phase(~p,~p,~p).~n",
              [StartPhase, StartType, Args]).

stop(_Data) ->
    ok.
\end{lstlisting}

\begin{notebox}[Start Phases 执行顺序]
Start Phases 按照 \texttt{start\_phases} 列表中定义的顺序依次执行。对于包含应用（Included Applications），子应用的 Start Phases 必须是父应用 Start Phases 的子集。
\end{notebox}


% ============================================================
% Section 8: Included Applications
% ============================================================
\section{包含应用（Included Applications）}

一个 Application 可以包含其他 Application。被包含的 Application 由包含它的 Application 负责启动和停止。

\subsection{配置 Included Applications}

\begin{lstlisting}[style=erlangstyle, caption={top\_app.app --- 包含 bsc 的 Application}]
{application, top_app,
 [{description, "Included Application Example"},
  {vsn, "1.0"},
  {modules, [top_app]},
  {applications, [kernel, stdlib, sasl]},
  {included_applications, [bsc]},
  {start_phases, [{start, []},
                  {admin, []},
                  {stop, []}]},
  {mod, {application_starter, [top_app, []]}}
 ]}.
\end{lstlisting}

\subsection{包含应用的回调模块}

\begin{lstlisting}[style=erlangstyle, caption={top\_app.erl --- 包含应用的回调模块}]
-module(top_app).
-behaviour(application).
-export([start/2, start_phase/3, stop/1]).

start(_Type, _Args) ->
    {ok, _Pid} = bsc_sup:start_link().

start_phase(StartPhase, StartType, Args) ->
    io:format("top_app:start_phase(~p,~p,~p).~n",
              [StartPhase, StartType, Args]).

stop(_Data) ->
    ok.
\end{lstlisting}

\begin{warnbox}[注意]
使用 \texttt{included\_applications} 时，\texttt{mod} 字段通常设置为 \texttt{\{application\_starter, [TopApp, []]\}}，这是 OTP 提供的一个辅助模块，用于协调包含应用的启动。被包含的 Application 不应该出现在 \texttt{applications} 依赖列表中。
\end{warnbox}


% ============================================================
% Section 9: Configuration
% ============================================================
\section{配置管理}

\subsection{sys.config 文件}

Application 的环境变量可以通过 \texttt{sys.config} 文件（或 \texttt{.config} 文件）在启动时配置：

\begin{lstlisting}[style=erlangstyle, caption={bsc.config --- 系统配置文件}]
[{sasl, [{errlog_type, all},
         {sasl_error_logger, tty}]},
 {bsc,  [{frequencies, [1,2,3,4,5,6]}]}].
\end{lstlisting}

\subsection{使用配置文件启动}

\begin{lstlisting}[style=shellstyle, caption={使用配置文件启动 Erlang 节点}]
$ erl -config bsc.config
\end{lstlisting}

\subsection{运行时读取配置}

\begin{lstlisting}[style=erlangstyle, caption={运行时读取和设置环境变量}]
%% Read configuration
application:get_env(bsc, frequencies).
%% => {ok, [1,2,3,4,5,6]}

%% Read with default value
application:get_env(bsc, max_retries, 3).
%% => 3

%% Set configuration at runtime
application:set_env(bsc, frequencies, [1,2,3,4,5,6,7,8]).
\end{lstlisting}

\begin{warnbox}[运行时修改配置的注意事项]
\texttt{application:set\_env/3,4} 应谨慎使用。配置参数的读取频率和时机取决于具体的 Application 实现。不当使用可能导致 Application 处于不一致的状态。

如果在 Application 加载之前调用 \texttt{set\_env}，\texttt{.app} 文件中定义的值会覆盖之前设置的值。可以使用 \texttt{persistent} 选项来避免这种情况。
\end{warnbox}


% ============================================================
% Section 10: Practical Demo
% ============================================================
\section{实验演示}

\subsection{启动 Application}

\begin{lstlisting}[style=shellstyle, caption={启动 BSC Application}]
$ erl -pa ebin -config bsc.config
Erlang/OTP 26 [erts-14.2.5.12] [source] [64-bit]

1> application:start(sasl).
ok
2> application:start(bsc).
ok
3> application:which_applications().
[{bsc,"Base Station Controller","1.0"},
 {sasl,"SASL  CXC 138 11","4.2.2"},
 {stdlib,"ERTS  CXC 138 10","5.2.3"},
 {kernel,"ERTS  CXC 138 10","9.2.4"}]
\end{lstlisting}

\subsection{使用 ensure\_all\_started}

\begin{lstlisting}[style=shellstyle, caption={使用 ensure\_all\_started 自动启动依赖}]
$ erl -pa ebin -config bsc.config
Erlang/OTP 26 [erts-14.2.5.12] [source] [64-bit]

1> application:ensure_all_started(bsc).
{ok,[sasl,bsc]}
\end{lstlisting}

\begin{tipbox}[ensure\_all\_started vs start]
\texttt{ensure\_all\_started/1} 会自动递归启动所有依赖的 Application，而 \texttt{start/1} 只启动指定的 Application，如果依赖未启动会返回错误。在实际项目中，推荐使用 \texttt{ensure\_all\_started/1}。
\end{tipbox}

\subsection{停止 Application}

\begin{lstlisting}[style=shellstyle, caption={停止 BSC Application}]
4> application:stop(bsc).
ok
=INFO REPORT==== ...
    application: bsc
    exited: stopped
    type: temporary
5> application:which_applications().
[{sasl,"SASL  CXC 138 11","4.2.2"},
 {stdlib,"ERTS  CXC 138 10","5.2.3"},
 {kernel,"ERTS  CXC 138 10","9.2.4"}]
\end{lstlisting}

\subsection{查询 Application 信息}

\begin{lstlisting}[style=shellstyle, caption={查询 Application 信息}]
6> application:start(bsc).
ok
7> application:get_env(bsc, frequencies).
{ok,[1,2,3,4,5,6]}
8> application:get_key(bsc, vsn).
{ok,"1.0"}
9> application:get_key(bsc, modules).
{ok,[bsc,bsc_sup,frequency,freq_overload,
     logger,simple_phone_sup,phone_fsm]}
10> application:get_supervisor(bsc).
{ok,<0.95.0>}
11> application:loaded_applications().
[{bsc,"Base Station Controller","1.0"},
 {sasl,"SASL  CXC 138 11","4.2.2"},
 {stdlib,"ERTS  CXC 138 10","5.2.3"},
 {kernel,"ERTS  CXC 138 10","9.2.4"}]
\end{lstlisting}


% ============================================================
% Section 11: Directory Structure
% ============================================================
\section{Application 目录结构}

标准的 OTP Application 目录结构如下：

\begin{lstlisting}[style=erlangstyle, caption={标准 OTP Application 目录结构}]
myapp-1.0/
  ebin/
    myapp.app          %% Application resource file
    myapp.beam         %% Compiled callback module
    myapp_sup.beam     %% Compiled supervisor
    ...
  src/
    myapp.erl          %% Application callback module
    myapp_sup.erl      %% Top supervisor
    ...
  include/             %% Header files (.hrl)
  priv/                %% Private data files
  doc/                 %% Documentation
\end{lstlisting}

使用 \texttt{rebar3} 构建工具时，目录结构略有不同：

\begin{lstlisting}[style=erlangstyle, caption={rebar3 项目目录结构}]
myapp/
  src/
    myapp.app.src      %% Application resource template
    myapp_app.erl      %% Application callback module
    myapp_sup.erl      %% Top supervisor
    ...
  include/
  priv/
  config/
    sys.config         %% System configuration
    vm.args            %% VM arguments
  rebar.config         %% Build configuration
\end{lstlisting}

\begin{notebox}[rebar3 的 .app.src]
使用 \texttt{rebar3} 时，\texttt{.app.src} 文件是模板，\texttt{rebar3} 会在编译时自动生成 \texttt{ebin/myapp.app} 文件，并自动填充 \texttt{modules} 列表。
\end{notebox}


% ============================================================
% Section 12: Common Pitfalls
% ============================================================
\section{常见陷阱}

\begin{itemize}[leftmargin=2em]
    \item \textbf{忘记声明依赖}：\texttt{applications} 列表中必须包含所有运行时依赖的 Application。遗漏会导致 \texttt{\{error, \{not\_started, App\}\}} 错误

    \item \textbf{循环依赖}：Application 之间不能有循环依赖。如果 A 依赖 B，B 又依赖 A，系统无法启动

    \item \textbf{start/2 返回值错误}：\texttt{start/2} 必须返回 \texttt{\{ok, Pid\}} 或 \texttt{\{ok, Pid, State\}}，其中 \texttt{Pid} 是 Top Supervisor 的 PID。返回其他值会导致启动失败

    \item \textbf{stop/1 中的清理}：\texttt{stop/1} 在所有进程终止\textbf{之后}才被调用。如果需要在进程终止之前执行清理，应使用 \texttt{prep\_stop/1}

    \item \textbf{环境变量的加载时机}：\texttt{application:set\_env/3} 在 Application 加载时会被 \texttt{.app} 文件中的值覆盖。使用 \texttt{\{persistent, true\}} 选项可以避免

    \item \textbf{included\_applications 的限制}：被包含的 Application 不能被独立启动，也不能出现在其他 Application 的 \texttt{applications} 依赖列表中

    \item \textbf{start\_phase 未定义}：如果 \texttt{.app} 文件中定义了 \texttt{start\_phases}，但回调模块没有导出 \texttt{start\_phase/3}，会导致 \texttt{undef} 错误
\end{itemize}


\end{document}
